\section{Учёт возможностей пользователя}
Цель данной работы подразумевает, что рекомендации рецептов основаны не только на предпочтениях пользователя, но и на его возможностях.
Следовательно, необходимо фильтровать список рекомендованных пользователю рецептов.
Критерий, определяющий, стоит ли рекомендовать рецепт, заключается в следующем: количество ингредиентов этого рецепта, недоступных пользователю, должно быть не больше $N$.
Далее описан алгоритм учёта возможностей пользователя.
Вводятся следующие обозначения:
\begin{itemize}
    \item $A$ -- множество ингредиентов, которые у пользователя в наличии;
    \item $R$ -- множество рецептов;
    \item $I(r), r \in R$ -- множество ингредиентов рецепта $r$.
\end{itemize}
Условие соответствия рецепта критерию выбора имеет вид~\eqref{eq:select_cond}:
\begin{equation}
	\label{eq:select_cond}
	|\{x: x \in I(r), x \notin A\}| \leqslant N.
\end{equation}
Алгоритм фильтрации $R$ содержит следующие шаги:
\begin{itemize}
    \item создать пустой список $L$;
    \item перебрать все рецепты $r \in R$;
    \item если условие~\ref{eq:select_cond} выполнено, то добавить $r$ в список $L$;
    \item конец цикла по $r$.
\end{itemize}
В результате получается список рецептов, которые рекомендованы пользователю.
При необходимости критерий~\eqref{eq:select_cond} дополняется проверкой оценки и тегов, присвоенных рецепту.